\section{What version 1.0.0 does}
\label{sec:functionality}

\subsection{Inference}

Table~\ref{tab:analyses} lists the analyses available, generated from
\code{registered_cells()} so that it matches the release it describes. Seven
cells are implemented. Six correspond to PhyloNet commands and are
reimplementations of published methods; the seventh, maximum-parsimony
inference under polyploidy, is a new implementation of the method PhyloNet
exposes as \code{InferNetwork_MP_Allopp} and differs from it in its proposal
kernel.

\begin{table}[t]
  \centering
  \scriptsize
  % Generated by benchmarks/make_analyses_table.py from phynetpy.infer.registered_cells().
% Do not edit by hand.
\setlength{\tabcolsep}{3pt}
\begin{tabular}{@{}l>{\raggedright\arraybackslash}p{15mm}>{\raggedright\arraybackslash}p{18mm}>{\raggedright\arraybackslash}p{16mm}>{\raggedright\arraybackslash}p{37mm}l>{\raggedright\arraybackslash}p{17mm}@{}}
\toprule
Engine & Data & Model & Criterion & Problem addressed & PhyloNet command & Status \\
\midrule
\texttt{MCMC\_GT} & Gene trees & MSC & Bayesian & Posterior over networks given gene trees \cite{Wen2016MCMCGT} & \texttt{MCMC\_GT} & Reimpl. \\
\texttt{InferNetwork\_ML} & Gene trees & MSC & Likelihood & Network topology and parameters by maximum likelihood \cite{Yu2014ML} & \texttt{InferNetwork\_ML} & Reimpl. \\
\texttt{InferNetwork\_MPL} & Gene trees & MSC & Pseudo-lik. & Network topology from gene-tree triplet frequencies \cite{Yu2015MPL} & \texttt{InferNetwork\_MPL} & Reimpl. \\
\texttt{MCMC\_SEQ} & Alignment & MSC & Bayesian & Joint posterior over networks and gene trees from sequences \cite{Wen2018MCMCSEQ} & \texttt{MCMC\_SEQ} & Reimpl. \\
\texttt{MCMC\_BiMarkers} & Markers & MSC & Bayesian & Posterior over networks from biallelic markers \cite{Zhu2018BiMarkers} & \texttt{MCMC\_BiMarkers} & Reimpl. \\
\texttt{MLE\_BiMarkers} & Markers & MSC & Likelihood & Exact marker likelihood of a given network \cite{Bryant2012SNAPP,Zhu2018BiMarkers} & \texttt{MLE\_BiMarkers} & Reimpl., scoring only \\
\texttt{MP\_Allop} & Gene trees & Allopolyploid & Parsimony & Parsimonious network under polyploidy, minimising deep coalescences \cite{Yan2022Allop,Than2009MDC} & \texttt{InferNetwork\_MP\_Allopp} & New impl. \\
\midrule
\multicolumn{7}{@{}l}{\emph{Cells that are well defined but carry no engine}} \\
--- & Gene trees & MSC & Parsimony & \emph{raises} \texttt{NotImplementedError} & \texttt{InferNetwork\_MP} & Sched. 1.0.0 \\
--- & Markers & MSC & Pseudo-lik. & \emph{raises} \texttt{NotImplementedError} & \texttt{MLE\_BiMarkers -pseudo} & Sched. 1.0.0 \\
--- & Alignment & MSC & Likelihood & \emph{raises} \texttt{NotImplementedError} & \texttt{---} & Not planned \\
\bottomrule
\end{tabular}

  \caption{Analyses in \phynetpy\ 1.0.0. Rows are generated from the dispatch
  registry. ``Reimplementation'' means the method was published elsewhere and
  implemented here from its description; the \code{MP_Allop} row is marked as
  a new implementation because its search uses a single ploidy-preserving move
  rather than the six moves of the original. The lower block lists
  combinations that are well defined but carry no engine, and therefore raise
  \code{NotImplementedError} rather than failing obscurely.%
  }
  \label{tab:analyses}
\end{table}

Three points about that table matter more than its length.

The first is what is \emph{not} in it. Maximum parsimony from gene trees under
diploid hybridization -- PhyloNet's \code{InferNetwork_MP} -- is a legal cell
with no engine behind it. The only parsimony code in \phynetpy\ is defined on
multiply-labelled trees for the allopolyploid model, so the diploid case is
genuinely absent rather than partially present. The biallelic
pseudo-likelihood is absent for the same kind of reason: the marker code
computes the full likelihood only. Both are scheduled for
1.0.0.\footnote{[TODO: verify at 1.0.0. At the time of writing these two cells
raise \code{NotImplementedError}; the manuscript must be updated with measured
behaviour once they are built, or this claim removed.]}

The second is a limitation inside a row that is otherwise marked available.
\code{MLE_BiMarkers} can score a given network but cannot search: its
\code{infer} path raises unconditionally. A reader comparing feature lists
would otherwise reasonably assume that maximum-likelihood network inference
from markers is available.

The third concerns \code{Bayesian(objective=PseudoLikelihood())}, which is
refused. A triplet pseudo-likelihood is not a normalised probability of the
data, so using it as a Metropolis--Hastings target does not yield a calibrated
posterior. The combination is mechanically easy to allow and would produce
numbers; it is rejected because those numbers would not mean what a reader of
the output would take them to mean.

Search behaviour is controlled by presets rather than by combinations of
booleans. Four are provided: \code{default}, \code{fast}, \code{accurate}, and
\code{phylonet}. The last exists to reproduce PhyloNet's
optimise-everything-per-topology behaviour, which is what makes cross-checking
against it meaningful; individual flags override the preset.

\subsection{Network analysis}

The analysis layer is usable without running any inference, and is where a
researcher comparing methods or summarising results spends time.

Nine dissimilarity measures are implemented: the $\mu$ (path-multiplicity)
distance \cite{Cardona2009MuDistance}, hardwired and softwired cluster
distances, Robinson--Foulds over hardwired clusters, tripartition distance,
displayed-tree distance, average and weighted average path distance
\cite{TODO_APD}, and a pairwise leaf distance; \code{Network} additionally
provides a rooted triplet distance. The cluster, tripartition and
displayed-tree measures accept a \code{normalize} argument returning a value in
$[0,1]$.

A separate module compares networks by the reticulations they assert rather
than by overall topology, through reticulation tripartitions, a dissimilarity
built on them, and precision and recall against a reference network. Whether an
inferred network recovers the \emph{right hybridizations} is the question a
biologist usually has, and a whole-network distance answers it only indirectly.

Decomposition covers displayed trees with their probabilities, the dominant
tree, hardwired and softwired clusters, biconnected components (blobs) and the
tree of blobs, network level, bridges and articulation points, and induced
subnetworks by taxon set.

\subsection{Simulation}

\code{simulate} is the generative inverse of \code{infer} and returns a
data-axis object, so a recovery check is two calls with nothing in between:

\begin{lstlisting}
sim = simulate(MSC(theta=0.02), taxa=6, n=200)
recovered = infer(sim, criterion=PseudoLikelihood())
\end{lstlisting}

Gene genealogies are simulated backwards in time within network branches under
the multispecies network coalescent, with lineages routed at reticulations
according to inheritance probabilities. Sequence evolution runs a continuous-time
Markov chain down each gene tree under JC69, HKY85 or GTR, single- or
multi-locus. Biallelic markers are simulated site by site, sampling an
independent genealogy per site. Species trees come from Yule or constant-rate
birth--death processes; networks are produced by drawing a tree and grafting
reticulations onto it.

Two limits are worth stating. Simulation under \code{Allopolyploid} raises
\code{NotImplementedError} -- only \code{MSC} has simulators. And because
reticulations are grafted onto a separately drawn tree, the network simulator
is not a birth--death--hybridization process in which hybridization and
speciation occur jointly, so it should not be used to study the distribution of
network shapes such a process induces.

One asymmetry catches users. The simulators write branch lengths in expected
substitutions per site while the gene-tree criteria read coalescent units, so a
round trip recovers the topology but not the lengths. The branch-length unit
tags described in Section~\ref{sec:architecture} make this visible rather than
silent.

\subsection{Data, I/O, and diagnostics}

Extended Newick, NEXUS, FASTA and VCF are supported for both reading and
writing. The extended Newick reader detects dialects and converts between them,
which matters because no single interchange convention is universal across
network software. PHYLIP is not supported.

Four alphabets are built in (DNA, RNA, protein, codon). Eight nucleotide
substitution models are provided -- JC, K80, F81, HKY, TN93, K81, SYM, GTR --
each parameterised by the free parameters its literature defines rather than by
a full exchangeability vector.

For posterior analysis, \phynetpy\ computes effective sample size,
autocorrelation time, standard error of the mean, highest-posterior-density
intervals, the Geweke statistic \cite{Geweke1992}, and the Gelman--Rubin
$\hat{R}$ across chains \cite{GelmanRubin1992}. Traces are written in Tracer
log format \cite{Rambaut2018Tracer} and tree samples in NEXUS, and Tracer logs
can be read back, so existing tooling applies to \phynetpy\ output without
conversion scripts. \code{reticulation_sweep} fits a range of reticulation
counts and reports information criteria, which addresses the standing problem
that a more reticulate network almost always fits the data better.
